%% kapa: exam=1 round=28 year=2017 period=1 form=A subject=real_estate subject_ko=부동산학원론 seq=28 number=108
%% source: Q-Net 감정평가사 기출문제 (cst00309) artlSeq=5199603 file=extracted/1008/5199603/2017년 제28회 감정평가사 1차시험 문제지/1교시 A형.pdf
%% conversion: kordoc 4.24.51 -> tools/gen1.py (auto); review=proofread-sonnet; flags=none
\begin{question}{108}
부동산시장에 관한 설명으로 옳은 것은?
\begin{choices}
  \item 부동산시장은 부동산 재화와 서비스가 교환되는 매커니즘이기 때문에 유형의 부동산 거래는 허용되며, 무형의 이용과 관련한 권리는 제외된다.
  \item 일반적으로 부동산시장은 일반시장에 비해 거래비용이 많이 들고, 수요자와 공급자의 시장진출입이 제약을 받게 되어 완전경쟁시장이 된다.
  \item 부동산의 입지성으로 인해 소유자는 해당 부동산의 활용과 가격결정에 있어서 입지 독점권(location monopoly)을 가지며, 이것은 하위시장의 형성과 관련 있다.
  \item 정부가 제품의 품질이나 규격을 통제하는 건축기준은 양적규제의 예로 들 수 있다.
  \item 준강성 효율적 시장은 공표된 것이건 그렇지 않은 것이건 어떠한 정보도 이미 가치에 반영되어 있는 시장이다.
\end{choices}
\end{question}
\answer{3}
